\scanpage{75}
% archive sheet number: 74
% printed form
% letterhead
UNIVERSITE\\
DES SCIENCES ET TECHNIQUES \hfill Montpellier, le 1er décembre 1977\\
\uline{DU LANGUEDOC}\\
JP/MD

% slanted
\margin{\handwritten{\unsure{88|BB}}}
Diffusion : Liste A\\
n\textdegree{} 6828

AVIS DE SOUTENANCE D'UNE THESE DE DOCTORAT DE SPECIALITE

SCIENCES DE LA TERRE - Mention GEOLOGIE APPLIQUEE

\uline{Nom et prénom du candidat} : ERRE Henry

\uline{Titre de la thèse} : Contribution à l'étude de l'hydrogéologie des Corbières orientales karstiques et pseudo-karstiques, et des émergences littorales des côtes calcaires du Languedoc-Roussillon.

\uline{Date, heure et lieu de soutenance prévus} : Mardi 13 décembre 1977 à 16 h 30, salle de cours 10 (Administration).

\uline{Nom du Directeur de thèse} : Monsieur le Professeur J. AVIAS

\uline{Désignation du laboratoire où la thèse est préparée} : Centre d'Etudes et de Recherches Géologiques et Hydrologiques.

\uline{Analyse de l'ouvrage} : L'auteur étudie les terrains jura-cretacés et triasiques en tant que systèmes aquifères. Ils le conduisent à examiner plus généralement les émergences littorales des côtes calcaires et à proposer une interprétation des anomalies thermiques et chimiques observées.

\uline{Durée approximative des recherches ayant abouti à la thèse} : 2 ans 1/2

\uline{Composition du jury proposé} : M. AVIAS - Président\\
\hspace*{7em}M. BOCQUILLON\\
\hspace*{7em}M. PLEGAT\\
\hspace*{7em}M. PALOC\\
\hspace*{7em}M. SALVAYRE

Le Secrétaire Général,\\
\handwritten{\unsure{E. Siau|E. Si}}\\
\uline{E. SIAU}
